\chapter{\texorpdfstring{Miscellany}{8. Miscellany}}

\section{\texorpdfstring{Contest Utilities}{8.1 Contest Utilities}}
\setcounter{section}{1}
\label{sec:8.1}
Small contest convenience helpers that are useful across many algorithms. These snippets avoid algorithm-specific policy and are meant to be pasted near the top of a solution file.

A personal template may also collect shortened versions of utilities elsewhere in the anthology, such as fast I/O and debugging, custom hashing, coordinate compression, bit operations, modular arithmetic, binary search, and GNU policy-based data structures. They remain in their dedicated sections here to avoid duplication and let each template be assembled to taste.

\begin{compactitem}
  \item \inlinecode{Rep(i,N)}, \inlinecode{For(i,L,H)}, \inlinecode{Rev(i,N)}, \inlinecode{Dwn(i,H,L)}, and \inlinecode{Each(x,C)} are compact loop macros.
  \item \inlinecode{All(c)} and \inlinecode{Rall(c)} expand to the forward or reverse iterator pair of a container or array.
  \item \inlinecode{sz(c)} returns the signed size of a container.
  \item \inlinecode{ckmin(a, b)} assigns \inlinecode{a = b} only when \inlinecode{b} \textless{} \inlinecode{a}, and returns whether the assignment happened.
  \item \inlinecode{ckmax(a, b)} assigns \inlinecode{a = b} only when \inlinecode{a} \textless{} \inlinecode{b}, and returns whether the assignment happened.
  \item \inlinecode{floor\_div(a, b)} and \inlinecode{ceil\_div(a, b)} divide signed integers with mathematical rounding toward negative or positive infinity. Requires nonzero \inlinecode{b} and the resulting quotient to be representable by the result type; in particular, the minimum value cannot be divided by $-1$.
  \item \inlinecode{make\_vec(n1, ..., nk, v)} returns a \inlinecode{k}-dimensional nested vector with the given sizes and every element initialized to \inlinecode{v}, whose type is deduced from \inlinecode{v}. The last argument is always the fill value, so \inlinecode{make\_vec(3, 5)} is a vector of three fives rather than a $3$ by $5$ grid.
  \item \inlinecode{indices(n)} returns the vector $[0, 1, \ldots, {\inlinecodemath{n}} - 1]$.
  \item \inlinecode{lb(a, x, comp = std::less\textless{}\textgreater{})} and \inlinecode{ub(a, x, comp = std::less\textless{}\textgreater{})} return lower and upper bound indices in random-access sequence \inlinecode{a}, which must be sorted according to \inlinecode{comp}.
  \item \inlinecode{sort\_unique(v, comp = std::less\textless{}\textgreater{})} sorts a vector and removes duplicates, which are the elements that \inlinecode{comp} deems equivalent.
  \item \inlinecode{argsort(a, comp = std::less\textless{}\textgreater{})} returns the indices of \inlinecode{a} ordered by their values according to \inlinecode{comp}; the order of indices whose values compare equal is unspecified.
  \item \inlinecode{find\_ptr(m, k)} returns a pointer to the value that key \inlinecode{k} maps to in the associative container \inlinecode{m}, or \inlinecode{nullptr} if the key is absent. The pointer is const only when \inlinecode{m} is, so \inlinecode{if (auto *p = find\_ptr(m, k))} both tests for the key and updates it in a single lookup.
  \item \inlinecode{get\_or(m, k, def = \{\})} returns the value that key \inlinecode{k} maps to in the associative container \inlinecode{m}, or \inlinecode{def} if the key is absent. Unlike \inlinecode{operator[]}, it never inserts and works on a const \inlinecode{m}.
  \item \inlinecode{erase\_one(c, x)} erases one copy of \inlinecode{x} from associative container \inlinecode{c}, asserting that it is present.
  \item \inlinecode{min\_heap\textless{}T\textgreater{}} is a min-heap alias.
  \item \inlinecode{y\_combinator(f)} wraps a recursive lambda so that it can call itself as its first argument.
\end{compactitem}

\begin{lstlisting}[style=cpp]
#include <algorithm>
#include <cassert>
#include <functional>
#include <iterator>
#include <numeric>
#include <queue>
#include <type_traits>
#include <utility>
#include <vector>

#define Rep(i, N)    for (int i = 0, _##i = (N); i < _##i; ++i)     //   0 to N-1
#define For(i, L, H) for (int i = (L), _##i = (H); i <= _##i; ++i)  //   L to H
#define Rev(i, N)    for (int i = (N); --i >= 0;)                   // N-1 to 0
#define Dwn(i, H, L) for (int i = (H), _##i = (L); i >= _##i; --i)  //   H to L
#define Each(x, C)   for (auto &x : (C))
#define All(C)       std::begin(C), std::end(C)
#define Rall(C)      std::rbegin(C), std::rend(C)

template<typename C> int sz(const C &c) { return static_cast<int>(c.size()); }
template<typename T, typename U> bool ckmin(T &a, const U &b) { return b < a && (a = b, true); }
template<typename T, typename U> bool ckmax(T &a, const U &b) { return a < b && (a = b, true); }

template<typename T>
T floor_div(T a, T b) {
  assert(b != 0);
  T q = a / b, r = a % b;
  return q - (r != 0 && ((r < 0) != (b < 0)));
}

template<typename T>
T ceil_div(T a, T b) {
  assert(b != 0);
  T q = a / b, r = a % b;
  return q + (r != 0 && ((r < 0) == (b < 0)));
}

template<typename T>
std::vector<T> make_vec(int n, const T &v) {
  return std::vector<T>(n, v);
}

template<typename... Args>
auto make_vec(int n, Args... args) {
  auto inner = make_vec(args...);
  return std::vector<decltype(inner)>(n, inner);
}

std::vector<int> indices(int n) {
  std::vector<int> v(n);
  std::iota(v.begin(), v.end(), 0);
  return v;
}

template<typename Seq, typename T, typename Compare = std::less<>>
int lb(const Seq &a, const T &x, Compare comp = Compare{}) {
  return static_cast<int>(std::lower_bound(a.begin(), a.end(), x, comp) - a.begin());
}

template<typename Seq, typename T, typename Compare = std::less<>>
int ub(const Seq &a, const T &x, Compare comp = Compare{}) {
  return static_cast<int>(std::upper_bound(a.begin(), a.end(), x, comp) - a.begin());
}

template<typename T, typename Compare = std::less<>>
void sort_unique(std::vector<T> &v, Compare comp = Compare{}) {
  std::sort(v.begin(), v.end(), comp);
  auto same = [&](const T &a, const T &b) { return !comp(a, b) && !comp(b, a); };
  v.erase(std::unique(v.begin(), v.end(), same), v.end());
}

template<typename Seq, typename Compare = std::less<>>
std::vector<int> argsort(const Seq &a, Compare comp = Compare{}) {
  std::vector<int> p = indices(static_cast<int>(a.size()));
  std::sort(p.begin(), p.end(), [&](int i, int j) { return comp(a[i], a[j]); });
  return p;
}

template<typename M>
auto find_ptr(M &m, const typename M::key_type &k) {
  auto it = m.find(k);
  return it == m.end() ? nullptr : &it->second;
}

template<typename M>
typename M::mapped_type get_or(
    const M &m, const typename M::key_type &k, const typename M::mapped_type &def = {}
) {
  auto it = m.find(k);
  return it == m.end() ? def : it->second;
}

template<typename C, typename T>
void erase_one(C &c, const T &x) {
  auto it = c.find(x);
  assert(it != c.end());
  c.erase(it);
}

template<typename T>
using min_heap = std::priority_queue<T, std::vector<T>, std::greater<T>>;

template<typename Fn>
class y_combinator_result {
  Fn fn;

 public:
  template<typename T>
  explicit y_combinator_result(T &&f) : fn(std::forward<T>(f)) {}

  template<typename... Args>
  decltype(auto) operator()(Args &&...args) {
    return fn(std::ref(*this), std::forward<Args>(args)...);
  }
};

template<typename Fn>
decltype(auto) y_combinator(Fn &&f) {
  return y_combinator_result<std::decay_t<Fn>>(std::forward<Fn>(f));
}
\end{lstlisting}
\Needspace*{4\baselineskip}
\textbf{Example Usage}\par
\begin{lstlisting}[style=cpp]
#include <map>
#include <set>
using namespace std;

int main() {
  int tot = 0;
  Rep(i, 4) tot += i;     // 0 + 1 + 2 + 3 = 6
  For(i, 2, 5) tot += i;  // 6 + (2 + 3 + 4 + 5) = 20
  Rev(i, 3) tot += i;     // 20 + (2 + 1 + 0) = 23
  Dwn(i, 8, 6) tot += i;  // 23 + (8 + 7 + 6) = 44
  assert(tot == 44);

  int x = 10;
  assert(ckmin(x, 7) && x == 7);
  assert(!ckmin(x, 9) && x == 7);
  assert(ckmax(x, 11) && x == 11);

  vector<int> v{1, 2, 3};
  assert(sz(v) == 3);
  sort(Rall(v));
  assert((v == vector<int>{3, 2, 1}));
  assert(lb(v, 2, greater<>{}) == 1 && ub(v, 2, greater<>{}) == 2);
  sort(All(v));
  assert(lb(v, 2) == 1 && ub(v, 2) == 2);
  assert(floor_div(7, -3) == -3);
  assert(floor_div(-7, 3) == -3);
  assert(ceil_div(7, -3) == -2);
  assert(ceil_div(-7, 3) == -2);

  auto grid = make_vec(2, 3, -1);
  grid[0][0] = 7;
  assert(sz(grid) == 2 && sz(grid[1]) == 3 && grid[1][0] == -1);
  v = {3, 1, 3, 2, 1};
  sort_unique(v);
  assert((v == vector<int>{1, 2, 3}));
  v = {3, 1, 3, 2, 1};
  sort_unique(v, greater<>{});
  assert((v == vector<int>{3, 2, 1}));
  assert((indices(4) == vector<int>{0, 1, 2, 3}));
  assert((argsort(vector<int>{30, 10, 20}) == vector<int>{1, 2, 0}));
  assert((argsort(vector<int>{30, 10, 20}, greater<>{}) == vector<int>{0, 2, 1}));

  map<int, int> cnt{{2, 5}};
  if (int *p = find_ptr(cnt, 2)) {
    *p += 10;  // A single lookup both tests for the key and updates it.
  }
  const map<int, int> &view = cnt;  // A const container yields a pointer to const.
  assert(*find_ptr(view, 2) == 15 && find_ptr(view, 7) == nullptr);
  assert(get_or(cnt, 2) == 15 && get_or(cnt, 7, -1) == -1 && sz(cnt) == 1);
  min_heap<int> pq;
  pq.push(3);
  pq.push(1);
  assert(pq.top() == 1);

  multiset<int> ms{1, 2, 2, 3};
  erase_one(ms, 2);
  assert(ms.count(2) == 1);

  // Recursive lambda with y_combinator.
  auto fib =
      y_combinator([](auto fib, int n) -> int { return n < 2 ? n : fib(n - 1) + fib(n - 2); });
  assert(fib(10) == 55);

  // Recursive lambda without y_combinator: pass the lambda to itself as the first argument.
  vector<vector<int>> adj{{1, 2}, {0}, {0}};
  int seen = 0;
  auto dfs = [&](auto &&dfs, int u, int p) -> void {
    seen++;
    Each(v, adj[u]) if (v != p) dfs(dfs, v, u);
  };
  dfs(dfs, 0, -1);
  assert(seen == 3);

  // Recursive lambdas are automatically supported in C++23.
#if __cplusplus >= 202302L
  auto fact = [&](this auto fact, int n) -> int { return n ? n * fact(n - 1) : 1; };
  assert(fact(5) == 120);
#endif
  return 0;
}
\end{lstlisting}

\section{\texorpdfstring{Fast IO}{8.2 Fast IO}}
\setcounter{section}{2}
\label{sec:8.2}
Fast input and output helpers. For most cases, it is enough to speed up standard I/O simply with \inlinecode{ios::sync\_with\_stdio(false); cin.tie(nullptr);}. The following classes only come into play when input is genuinely huge or \inlinecode{iostream} overhead is measurable.

\begin{compactitem}
  \item \inlinecode{set\_in(name)}, \inlinecode{set\_out(name)}, and \inlinecode{set\_io(iname, oname)} redirect standard input/output to files. For example, \inlinecode{set\_io("task.in", "task.out")} is convenient for USACO-style problems. An empty name leaves that stream untouched, so \inlinecode{set\_io("", "")} is a no-op that can stay in place when the same solution is submitted to a judge that pipes its input.
\end{compactitem}

File-redirection failures throw \inlinecode{std::runtime\_error}.

\begin{lstlisting}[style=cpp]
#include <cctype>
#include <cstddef>
#include <cstdio>
#include <cstdlib>
#include <limits>
#include <stdexcept>
#include <string>
#include <type_traits>

void set_in(const std::string &name) {
  if (!name.empty() && !freopen(name.c_str(), "r", stdin)) {
    throw std::runtime_error("Failed to open input file: " + name);
  }
}

void set_out(const std::string &name) {
  if (!name.empty() && !freopen(name.c_str(), "w", stdout)) {
    throw std::runtime_error("Failed to open output file: " + name);
  }
}

void set_io(const std::string &iname, const std::string &oname) {
  set_in(iname);
  set_out(oname);
}
\end{lstlisting}
When only integers are read, one loop over a character reader is the whole technique and is short enough to copy or retype: skip anything that starts no number, then fold digits in with \inlinecode{x * 10 + d}. Which character reader it uses decides everything, since the per-character mutex that the locked \inlinecode{getchar()} takes costs more than the parsing does. The example below times the alternatives against each other on the same input.

\begin{compactitem}
  \item \inlinecode{read\_int(x)} reads one integer from \inlinecode{stdin} into \inlinecode{x}, skipping any leading character that cannot start a number. A leading \inlinecode{-} is honored, so \inlinecode{Int} must be signed for negative input. Behavior is undefined at end of file.
\end{compactitem}

\begin{lstlisting}[style=cpp]
// The unlocked reader is where nearly all of the speed comes from. It is POSIX, not standard C++.
inline int read_char() {
#if defined(_WIN32)
  return _getchar_nolock();
#elif defined(__unix__) || defined(__APPLE__)
  return getchar_unlocked();
#else
  return getchar();
#endif
}

template<typename Int>
void read_int(Int &x) {
  int c, s = 1;
  while ((c = read_char()) != '-' && (c < '0' || c > '9'));
  if (c == '-') s = -1, c = read_char();
  for (x = 0; '0' <= c && c <= '9'; c = read_char()) x = x * 10 + c - '0';
  x *= s;
}
\end{lstlisting}
\inlinecode{FastInput} reads large byte blocks with \inlinecode{fread()} and parses tokens from an internal buffer.

\begin{compactitem}
  \item \inlinecode{FastInput in(file = stdin)} constructs a faster reader from a \inlinecode{FILE*}.
  \item \inlinecode{in \textgreater{}\textgreater{} x} reads a non-whitespace token into \inlinecode{char}, \inlinecode{std::string}, integral types, or floating-point types.
\end{compactitem}

Detected read failures throw \inlinecode{std::runtime\_error}. End-of-file is not an error, and the parser otherwise assumes valid input. Floating-point input is provided for convenience, not as the main performance path.

\begin{lstlisting}[style=cpp]
class FastInput {
  static constexpr int BUF_SIZE = 1 << 20;
  FILE *file;
  char buf[BUF_SIZE];
  int pos = 0, len = 0;

  char get_char() {
    if (pos == len) {
      len = static_cast<int>(fread(buf, 1, BUF_SIZE, file));
      pos = 0;
      if (len == 0) {
        if (ferror(file)) {
          throw std::runtime_error("Failed to read input.");
        }
        return 0;
      }
    }
    return buf[pos++];
  }

  void skip_blanks() {
    while (true) {
      char c = get_char();
      if (!c) {
        return;
      }
      if (!std::isspace(static_cast<unsigned char>(c))) {
        --pos;
        return;
      }
    }
  }

 public:
  explicit FastInput(FILE *file_ = stdin) : file(file_) {}

  FastInput(const FastInput &) = delete;
  FastInput &operator=(const FastInput &) = delete;

  FastInput &operator>>(char &c) {
    skip_blanks();
    c = get_char();
    return *this;
  }

  FastInput &operator>>(std::string &s) {
    skip_blanks();
    s.clear();
    while (true) {
      char c = get_char();
      if (!c || std::isspace(static_cast<unsigned char>(c))) {
        break;
      }
      s += c;
    }
    return *this;
  }

  FastInput &operator>>(bool &x) {
    int val;
    *this >> val;
    x = val != 0;
    return *this;
  }

  template<typename T>
  typename std::enable_if<
      std::is_integral<T>::value && !std::is_same<T, bool>::value, FastInput &>::type
  operator>>(T &x) {
    skip_blanks();
    char c = get_char();
    bool neg = false;
    if (c == '-') {
      neg = true;
      c = get_char();
    }
    using U = typename std::make_unsigned<T>::type;
    U val = 0;
    while (c && !std::isspace(static_cast<unsigned char>(c))) {
      val = U{10} * val + U(c - '0');
      c = get_char();
    }
    if constexpr (std::is_signed<T>::value) {
      x = neg ? T(U{0} - val) : T(val);
    } else {
      x = T(val);
    }
    return *this;
  }

  template<typename T>
  typename std::enable_if<std::is_floating_point<T>::value, FastInput &>::type operator>>(T &x) {
    std::string s;
    *this >> s;
    x = static_cast<T>(std::strtold(s.c_str(), nullptr));
    return *this;
  }
};
\end{lstlisting}
\inlinecode{FastOutput} formats values into an internal buffer and writes them in large blocks with \inlinecode{fwrite()}.

\begin{compactitem}
  \item \inlinecode{FastOutput out(file = stdout)} constructs a faster writer to a \inlinecode{FILE*}.
  \item \inlinecode{out \textless{}\textless{} x} writes \inlinecode{char}, C strings, \inlinecode{std::string}, integral types, or floating-point types.
  \item \inlinecode{out.flush()} writes any buffered output immediately.
\end{compactitem}

Detected write failures throw \inlinecode{std::runtime\_error}. Destruction flushes on a best-effort basis; call \inlinecode{flush()} explicitly to detect a final write failure. Floating-point output is provided for convenience, not as the main performance path.

\begin{lstlisting}[style=cpp]
class FastOutput {
  static constexpr int BUF_SIZE = 1 << 20;
  FILE *file;
  char buf[BUF_SIZE];
  int pos = 0;

  bool flush_buf() {
    int n = pos;
    pos = 0;
    return fwrite(buf, 1, n, file) == static_cast<std::size_t>(n);
  }

  void put_char(char c) {
    if (pos == BUF_SIZE) {
      flush();
    }
    buf[pos++] = c;
  }

 public:
  explicit FastOutput(FILE *file_ = stdout) : file(file_) {}

  ~FastOutput() { flush_buf(); }
  FastOutput(const FastOutput &) = delete;
  FastOutput &operator=(const FastOutput &) = delete;

  void flush() {
    if (!flush_buf()) {
      throw std::runtime_error("Failed to write output.");
    }
  }

  FastOutput &operator<<(char c) {
    put_char(c);
    return *this;
  }

  FastOutput &operator<<(const char *s) {
    while (*s) {
      put_char(*s++);
    }
    return *this;
  }

  FastOutput &operator<<(const std::string &s) {
    for (char c : s) {
      put_char(c);
    }
    return *this;
  }

  FastOutput &operator<<(bool x) {
    put_char(x ? '1' : '0');
    return *this;
  }

  template<typename T>
  typename std::enable_if<
      std::is_integral<T>::value && !std::is_same<T, bool>::value, FastOutput &>::type
  operator<<(T x) {
    using U = typename std::make_unsigned<T>::type;
    U val = U(x);
    if constexpr (std::is_signed<T>::value) {
      if (x < 0) {
        put_char('-');
        val = U{0} - val;
      }
    }
    char s[std::numeric_limits<U>::digits10 + 1];
    int n = 0;
    do {
      s[n++] = char('0' + val % 10);
      val /= 10;
    } while (val);
    while (n--) {
      put_char(s[n]);
    }
    return *this;
  }

  template<typename T>
  typename std::enable_if<std::is_floating_point<T>::value, FastOutput &>::type operator<<(T x) {
    char s[64];
    std::snprintf(
        s, sizeof(s), "%.*Lg", std::numeric_limits<T>::max_digits10, static_cast<long double>(x)
    );
    return *this << s;
  }
};
\end{lstlisting}
\Needspace*{4\baselineskip}
\textbf{Example Usage}\par
\begin{lstlisting}[style=cpp]
#include <cassert>
#include <ctime>
#include <filesystem>
#include <fstream>
#include <iostream>
using namespace std;

int main() {
  // set_io("file.in", "file.out");

  FILE *input = tmpfile();
  fputs("42 hello 3.5 -2147483648 1 Z 1234567890123\n", input);
  rewind(input);

  FastInput in(input);
  int x;
  string s;
  double y;
  int z;
  bool flag;
  char ch;
  long long big;
  in >> x >> s >> y >> z >> flag >> ch >> big;
  assert(x == 42 && s == "hello" && y == 3.5 && z == -2147483648);
  assert(flag && ch == 'Z' && big == 1234567890123LL);
  fclose(input);

  FILE *output = tmpfile();
  {
    FastOutput out(output);
    out << x << ' ' << s << ' ' << y << ' ' << z << ' ' << flag << ' ' << ch << ' ' << big << '\n';
    out.flush();
  }
  rewind(output);
  char buf[96] = {};
  assert(fgets(buf, sizeof(buf), output));
  assert(string(buf) == "42 hello 3.5 -2147483648 1 Z 1234567890123\n");
  fclose(output);

  // Benchmark various ways of reading the same million integers from a scratch file.
  const int n = 1000000;
  string text;
  for (int i = 0; i < n; i++) {
    text += to_string(i) + ' ';
  }
  string path = (filesystem::temp_directory_path() / "fast_io.tmp").string();
  ofstream(path) << text;
  auto time_reads = [&](const char *name, auto read_one) {
    set_in(path);
    clock_t start = clock();
    for (int i = 0, v = 0; i < n; i++) read_one(v);
    printf("%-22s %6.3f s\n", name, double(clock() - start) / CLOCKS_PER_SEC);
  };
  printf("Reading %d integers:\n", n);
  time_reads("cin", [](int &v) { cin >> v; });
  time_reads("scanf", [](int &v) { scanf("%d", &v); });
  time_reads("read_int (unlocked)", [](int &v) { read_int(v); });
  time_reads("FastInput", [](int &v) {
    static FastInput in;
    in >> v;
  });

  auto time_writes = [&](const char *name, auto write_all) {
    clock_t start = clock();
    write_all(n);
    printf("%-22s %6.3f s\n", name, double(clock() - start) / CLOCKS_PER_SEC);
  };
  printf("\nWriting %d integers:\n", n);
  time_writes("ofstream", [&](int k) {
    ofstream out(path);
    for (int i = 0; i < k; i++) out << i << ' ';
  });
  time_writes("fprintf", [&](int k) {
    FILE *f = fopen(path.c_str(), "w");
    for (int i = 0; i < k; i++) fprintf(f, "%d ", i);
    fclose(f);
  });
  time_writes("FastOutput", [&](int k) {
    FILE *f = fopen(path.c_str(), "w");
    {
      FastOutput out(f);
      for (int i = 0; i < k; i++) out << i << ' ';
    }
    fclose(f);
  });
  remove(path.c_str());
  return 0;
}
\end{lstlisting}
\begin{exampleoutput}
Reading 1000000 integers:
cin                     0.959 s
scanf                   0.061 s
read_int (unlocked)     0.010 s
FastInput               0.021 s

Writing 1000000 integers:
ofstream                0.037 s
fprintf                 0.071 s
FastOutput              0.009 s
\end{exampleoutput}

\section{\texorpdfstring{Debugging}{8.3 Debugging}}
\setcounter{section}{3}
\label{sec:8.3}
Local-only debug printing for contest code. \inlinecode{dbg}, \inlinecode{pr}, \inlinecode{dbg\_rows}, and \inlinecode{pr\_rows} expand to no-ops unless \inlinecode{LOCAL} is defined (e.g. via the compiler flag \inlinecode{-DLOCAL}), so they can remain in submitted code without producing output.

\begin{compactitem}
  \item \inlinecode{dbg\_repr(x)} returns a string representation of arithmetic types, strings, pairs, tuples, iterable containers, and custom types that support insertion into \inlinecode{std::ostream}.
  \item \inlinecode{dbg(a, b, c)} prints the function name, line number, argument names, and values to \inlinecode{std::cerr}.
  \item \inlinecode{pr(a, b, c)} prints only the values to \inlinecode{std::cerr}; \inlinecode{pr()} prints a blank line.
  \item \inlinecode{dbg\_rows(rows, col\_width, limit = -1)} prints the function name, line number, argument name, and each iterable row on a separate indexed line. A positive \inlinecode{col\_width} right-aligns each element in a field of that width. Pass $0$ for ordinary container formatting. A nonnegative \inlinecode{limit} prints at most that many rows and elements per row, marking omitted entries with \inlinecode{...}.
  \item \inlinecode{pr\_rows(rows, col\_width, limit = -1)} prints only the indexed rows with the same formatting.
\end{compactitem}

Avoid naming these helpers \inlinecode{to\_string()} unless you want them mixed into overload resolution with \inlinecode{std::to\_string()} and other user-defined \inlinecode{to\_string()} functions. A separate name keeps debug formatting local to this snippet.

\begin{lstlisting}[style=cpp]
#include <cstddef>
#include <iomanip>
#include <iostream>
#include <sstream>
#include <string>
#include <tuple>
#include <type_traits>
#include <utility>

template<typename T, typename = void>
struct is_iterable : std::false_type {};

template<typename T>
struct is_iterable<
    T, std::void_t<decltype(std::begin(std::declval<T>())), decltype(std::end(std::declval<T>()))>>
    : std::true_type {};

template<typename T, typename = void>
struct is_streamable : std::false_type {};

template<typename T>
struct is_streamable<
    T, std::void_t<decltype(std::declval<std::ostream &>() << std::declval<const T &>())>>
    : std::true_type {};

std::string dbg_repr(const std::string &s) { return '"' + s + '"'; }
std::string dbg_repr(const char *s) { return dbg_repr(std::string(s)); }
std::string dbg_repr(char c) { return std::string("'") + c + "'"; }
std::string dbg_repr(bool b) { return b ? "true" : "false"; }

template<typename T>
typename std::enable_if<
    std::is_arithmetic<T>::value && !std::is_same<T, char>::value &&
        !std::is_same<T, bool>::value && !std::is_floating_point<T>::value,
    std::string>::type
dbg_repr(T x) {
  std::ostringstream out;
  out << x;
  return out.str();
}

template<typename T>
typename std::enable_if<std::is_floating_point<T>::value, std::string>::type dbg_repr(T x) {
  std::ostringstream out;
  out << std::setprecision(10) << x;  // Set floating point dbg precision here.
  return out.str();
}

template<typename T>
typename std::enable_if<
    is_streamable<T>::value && !is_iterable<T>::value && !std::is_arithmetic<T>::value &&
        !std::is_convertible<T, std::string>::value,
    std::string>::type
dbg_repr(const T &x) {
  std::ostringstream out;
  out << x;
  return out.str();
}

template<typename A, typename B>
std::string dbg_repr(const std::pair<A, B> &p) {
  return "(" + dbg_repr(p.first) + ", " + dbg_repr(p.second) + ")";
}

template<typename Tuple, std::size_t... Is>
std::string dbg_tuple_repr(const Tuple &t, std::index_sequence<Is...>) {
  std::string res = "(";
  bool first = true;
  ((res += (first ? (first = false, "") : ", ") + dbg_repr(std::get<Is>(t))), ...);
  return res + ")";
}

template<typename... Ts>
std::string dbg_repr(const std::tuple<Ts...> &t) {
  return dbg_tuple_repr(t, std::index_sequence_for<Ts...>{});
}

template<typename T>
typename std::enable_if<
    is_iterable<T>::value && !std::is_convertible<T, std::string>::value, std::string>::type
dbg_repr(const T &v) {
  std::string res = "{";
  bool first = true;
  for (const auto &x : v) {
    if (!first) {
      res += ", ";
    }
    first = false;
    res += dbg_repr(x);
  }
  return res + "}";
}

template<typename Head, typename... Tail>
void dbg_out(bool leading_space, const Head &head, const Tail &...tail) {
  if (leading_space) {
    std::cerr << ' ';
  }
  std::cerr << dbg_repr(head);
  ((std::cerr << ' ' << dbg_repr(tail)), ...);
  std::cerr << '\n';
}

void pr_out() {
  std::cerr << '\n';
}

template<typename Head, typename... Tail>
void pr_out(const Head &head, const Tail &...tail) {
  dbg_out(false, head, tail...);
}

template<typename Rows>
void dbg_rows_out(const Rows &rows, int col_width = 0, int limit = -1) {
  int i = 0;
  for (const auto &row : rows) {
    if (i == limit) {
      std::cerr << "...\n";
      break;
    }
    std::cerr << i++ << ':' << (col_width > 0 ? "" : " {");
    int j = 0;
    for (const auto &x : row) {
      if (j == limit) {
        std::cerr << (col_width > 0 ? " ..." : (j == 0 ? "..." : ", ..."));
        break;
      }
      std::cerr << (col_width > 0 ? " " : (j == 0 ? "" : ", ")) << std::setw(col_width)
                << dbg_repr(x);
      j++;
    }
    std::cerr << (col_width > 0 ? "\n" : "}\n");
  }
}

#ifdef LOCAL
#define dbg(...)                                                                  \
  std::cerr << "[" << __func__ << ":" << __LINE__ << "] " << #__VA_ARGS__ << ":", \
      dbg_out(true, __VA_ARGS__)
#define pr(...) pr_out(__VA_ARGS__)
#define dbg_rows(rows, ...)                                                  \
  std::cerr << "[" << __func__ << ":" << __LINE__ << "] " << #rows << ":\n", \
      dbg_rows_out(rows, __VA_ARGS__)
#define pr_rows(rows, ...) dbg_rows_out(rows, __VA_ARGS__)
#else
#define dbg(...) ((void)0)
#define pr(...) ((void)0)
#define dbg_rows(...) ((void)0)
#define pr_rows(...) ((void)0)
#endif
\end{lstlisting}
\Needspace*{4\baselineskip}
\textbf{Example Usage}\par
\begin{lstlisting}[style=cpp]
#include <cassert>
#include <map>
#include <vector>
using namespace std;

struct Point {
  int x, y;
};

// Can overload dbg_repr to dbg your custom structs.
string dbg_repr(const Point &p) {
  return "Point(" + dbg_repr(p.x) + ", " + dbg_repr(p.y) + ")";
}

struct Edge {
  int u, v, w;
};

string dbg_repr(const Edge &e) {
  return "Edge{" + dbg_repr(e.u) + "->" + dbg_repr(e.v) + ", w=" + dbg_repr(e.w) + "}";
}

struct Streamable {
  int value;

  friend ostream &operator<<(ostream &out, const Streamable &x) {
    return out << "Streamable(" << x.value << ")";
  }
};

int main() {
  vector<int> v{1, 2, 3};
  auto p = make_pair(4, string("x"));
  auto t = make_tuple(1, true, string("ok"));
  map<string, int> mp{{"a", 1}, {"b", 2}};
  Point pt{2, 3};
  Edge e{0, 1, 5};
  double d = 3.141592653589793;

  assert(dbg_repr(v) == "{1, 2, 3}");
  assert(dbg_repr(p) == "(4, \"x\")");
  assert(dbg_repr(t) == "(1, true, \"ok\")");
  assert(dbg_repr(mp) == "{(\"a\", 1), (\"b\", 2)}");
  assert(dbg_repr(pt) == "Point(2, 3)");
  assert(dbg_repr(e) == "Edge{0->1, w=5}");
  assert(dbg_repr(Streamable{6}) == "Streamable(6)");
  assert(dbg_repr(vector<Streamable>{{7}, {8}}) == "{Streamable(7), Streamable(8)}");
  assert(dbg_repr(1.0 / 3.0) == "0.3333333333");
  assert(dbg_repr(d) == "3.141592654");

  dbg(v);
  dbg(p, t, mp);
  dbg(pt, e);
  pr("plain output", d);

  vector<vector<int>> adj{{1, 2}, {0}, {0}};
  pr();
  dbg_rows(adj, 0);

  vector<vector<int>> grid{{1, 20, -3}, {400, 5, 6}, {7, 8, 9}};
  pr();
  dbg_rows(grid, 3, 2);
  return 0;
}
\end{lstlisting}
\begin{exampleoutput}
[main:250] v: {1, 2, 3}
[main:251] p, t, mp: (4, "x") (1, true, "ok") {("a", 1), ("b", 2)}
[main:252] pt, e: Point(2, 3) Edge{0->1, w=5}
"plain output" 3.141592654

[main:257] adj:
0: {1, 2}
1: {0}
2: {0}

[main:261] grid:
0:   1  20 ...
1: 400   5 ...
...
\end{exampleoutput}

\section{\texorpdfstring{Stress-Testing}{8.4 Stress-Testing}}
\setcounter{section}{4}
\label{sec:8.4}
Provides random input generators and a harness that compares an optimized solution against a trusted brute force implementation. Small tests keep brute force practical and produce manageable counterexamples. The generators create common contest inputs such as distinct samples, permutations, graphs, trees, and integer compositions for stress tests, randomized tests, and benchmarks.

\begin{compactitem}
  \item \inlinecode{rng} is a global 64-bit Mersenne Twister. Call \inlinecode{rng.seed(seed)} before a stress test to make it reproducible. Use an explicit seed so generated failures are reproducible.
  \item \inlinecode{rand\_int(lo, hi)} returns a uniformly random integer in range $[{\inlinecodemath{lo}}, {\inlinecodemath{hi}}]$.
  \item \inlinecode{rand\_biased(lo, hi, bias)} returns an integer in $[{\inlinecodemath{lo}}, {\inlinecodemath{hi}}]$. Positive \inlinecode{bias} favors larger values by taking the maximum of \inlinecode{bias + 1} samples, negative \inlinecode{bias} similarly favors smaller values, and zero is uniform.
  \item \inlinecode{rand\_real(lo = 0.0, hi = 1.0)} returns a uniformly random real number in the half-open range $[{\inlinecodemath{lo}}, {\inlinecodemath{hi}})$.
  \item \inlinecode{rand\_choice(values)} returns a uniformly random element of the nonempty vector \inlinecode{values}.
  \item \inlinecode{rand\_weighted(weights)} returns a sampled 0-based index with probability proportional to its nonnegative weight.
  \item \inlinecode{rand\_vec(n, lo, hi)} returns \inlinecode{n} independent uniformly random integers in $[{\inlinecodemath{lo}}, {\inlinecodemath{hi}}]$.
  \item \inlinecode{rand\_str(n, alphabet)} returns a length-\inlinecode{n} string whose characters are sampled independently and uniformly from the nonempty string \inlinecode{alphabet}.
  \item \inlinecode{rand\_perm(n, first = 0)} returns a random permutation of the \inlinecode{n} consecutive integers beginning at \inlinecode{first}.
  \item \inlinecode{rand\_distinct(k, lo, hi)} returns \inlinecode{k} distinct integers sampled uniformly from $[{\inlinecodemath{lo}}, {\inlinecodemath{hi}}]$.
  \item \inlinecode{rand\_tree(n)} returns a uniformly random labeled tree on nodes $[0, {\inlinecodemath{n}})$ as an edge list using a Prüfer code.
  \item \inlinecode{rand\_graph(n, m, directed = false, connected = false)} returns a random simple graph with \inlinecode{n} nodes numbered $[0, {\inlinecodemath{n}})$ and exactly \inlinecode{m} edges. Without connectivity it is uniform; with connectivity it starts from \inlinecode{rand\_tree(n)} and is not uniform over connected graphs. A connected directed graph is weakly connected.
  \item \inlinecode{rand\_composition(parts, sum, min\_part)} uniformly splits \inlinecode{sum} into \inlinecode{parts} ordered integers, each at least \inlinecode{min\_part}.
\end{compactitem}

\Needspace*{3\baselineskip}
\textbf{Time Complexity}:
\begin{compactitem}
  \item $O(1)$ per call to \inlinecode{rand\_int()}, \inlinecode{rand\_real()}, and \inlinecode{rand\_choice()}.
  \item $O(|b| + 1)$ per call to \inlinecode{rand\_biased()} where $b$ is the bias.
  \item $O(n)$ per call to \inlinecode{rand\_weighted()}, \inlinecode{rand\_vec()}, \inlinecode{rand\_str()}, and \inlinecode{rand\_perm()}, where $n$ is the number of weights or generated elements.
  \item $O(k)$ expected per call to \inlinecode{rand\_distinct()}.
  \item $O(n \log n)$ per call to \inlinecode{rand\_tree()}.
  \item $O(m)$ expected per call to \inlinecode{rand\_graph()} without connectivity, or $O(n \log n + m)$ with it.
  \item $O(k \log k)$ per call to \inlinecode{rand\_composition()}, where $k$ is \inlinecode{parts}.
\end{compactitem}

\Needspace*{3\baselineskip}
\textbf{Space Complexity}:
\begin{compactitem}
  \item $O(1)$ auxiliary for the scalar generators and \inlinecode{rand\_choice()}.
  \item $O(n)$ auxiliary for \inlinecode{rand\_weighted()}.
  \item $O(n)$ for the values returned by \inlinecode{rand\_vec()}, \inlinecode{rand\_str()}, \inlinecode{rand\_perm()}, and \inlinecode{rand\_tree()}.
  \item $O(k)$ auxiliary and output space for \inlinecode{rand\_distinct()} and \inlinecode{rand\_composition()}.
  \item $O(m)$ auxiliary and output space for \inlinecode{rand\_graph()}.
\end{compactitem}

\begin{lstlisting}[style=cpp]
#include <algorithm>
#include <cassert>
#include <climits>
#include <cmath>
#include <cstdint>
#include <functional>
#include <initializer_list>
#include <numeric>
#include <optional>
#include <queue>
#include <random>
#include <string>
#include <type_traits>
#include <unordered_set>
#include <utility>
#include <vector>

std::mt19937_64 rng;

template<typename Int>
Int rand_int(Int lo, Int hi) {
  static_assert(std::is_integral<Int>::value, "rand_int() requires an integral type");
  assert(lo <= hi);
  return std::uniform_int_distribution<Int>(lo, hi)(rng);
}

template<typename Int>
Int rand_biased(Int lo, Int hi, int bias) {
  Int value = rand_int(lo, hi);
  while (bias > 0) {
    value = std::max(value, rand_int(lo, hi));
    bias--;
  }
  while (bias < 0) {
    value = std::min(value, rand_int(lo, hi));
    bias++;
  }
  return value;
}

double rand_real(double lo = 0.0, double hi = 1.0) {
  assert(lo < hi);
  return std::uniform_real_distribution<double>(lo, hi)(rng);
}

template<typename T>
T rand_choice(const std::vector<T> &values) {
  assert(!values.empty() && values.size() <= INT_MAX);
  return values[rand_int(0, static_cast<int>(values.size()) - 1)];
}

template<typename W>
int rand_weighted(const std::vector<W> &weights) {
  static_assert(std::is_arithmetic<W>::value, "weights must be numeric");
  assert(!weights.empty() && weights.size() <= INT_MAX);
  std::vector<double> probabilities;
  probabilities.reserve(weights.size());
  double total = 0;
  for (const W &value : weights) {
    double weight = static_cast<double>(value);
    assert(std::isfinite(weight) && weight >= 0);
    probabilities.push_back(weight);
    total += weight;
  }
  assert(std::isfinite(total) && total > 0);
  return std::discrete_distribution<int>(probabilities.begin(), probabilities.end())(rng);
}

template<typename W>
int rand_weighted(std::initializer_list<W> weights) {
  return rand_weighted(std::vector<W>(weights));
}

template<typename Int>
std::vector<Int> rand_vec(int n, Int lo, Int hi) {
  assert(n >= 0 && lo <= hi);
  std::vector<Int> values(n);
  for (Int &x : values) {
    x = rand_int(lo, hi);
  }
  return values;
}

std::string rand_str(int n, const std::string &alphabet) {
  assert(n >= 0 && !alphabet.empty() && alphabet.size() <= INT_MAX);
  std::string s(n, ' ');
  for (char &c : s) {
    c = alphabet[rand_int(0, static_cast<int>(alphabet.size()) - 1)];
  }
  return s;
}

std::vector<int> rand_perm(int n, int first = 0) {
  assert(n >= 0 && (n == 0 || static_cast<int64_t>(first) + n - 1 <= INT_MAX));
  std::vector<int> p(n);
  std::iota(p.begin(), p.end(), first);
  std::shuffle(p.begin(), p.end(), rng);
  return p;
}

std::vector<int> rand_distinct(int k, int lo, int hi) {
  int64_t count = static_cast<int64_t>(hi) - lo + 1;
  assert(lo <= hi && 0 <= k && k <= count);
  std::unordered_set<int> used;
  used.reserve(k);
  std::vector<int> values;
  values.reserve(k);
  // Floyd's algorithm avoids materializing the entire interval.
  for (int64_t j = static_cast<int64_t>(hi) - k + 1; j <= hi; ++j) {
    int x = rand_int(lo, static_cast<int>(j));
    if (!used.insert(x).second) {
      x = static_cast<int>(j);
      used.insert(x);
    }
    values.push_back(x);
  }
  std::shuffle(values.begin(), values.end(), rng);
  return values;
}

std::vector<std::pair<int, int>> rand_tree(int n) {
  assert(n >= 1);
  if (n == 1) {
    return {};
  }
  std::vector<int> degree(n, 1), code(n - 2);
  for (int &u : code) {
    degree[u = rand_int(0, n - 1)]++;
  }
  std::priority_queue<int, std::vector<int>, std::greater<int>> leaves;
  for (int u = 0; u < n; ++u) {
    if (degree[u] == 1) {
      leaves.push(u);
    }
  }
  std::vector<std::pair<int, int>> edges;
  for (int u : code) {
    int leaf = leaves.top();
    leaves.pop();
    edges.emplace_back(leaf, u);
    if (--degree[u] == 1) {
      leaves.push(u);
    }
  }
  int u = leaves.top();
  leaves.pop();
  edges.emplace_back(u, leaves.top());
  std::shuffle(edges.begin(), edges.end(), rng);
  return edges;
}

std::vector<std::pair<int, int>> rand_graph(
    int n, int m, bool directed = false, bool connected = false
) {
  assert(n >= 0 && m >= 0 && (!connected || n >= 1));
  int64_t total =
      directed ? static_cast<int64_t>(n) * (n - 1) : static_cast<int64_t>(n) * (n - 1) / 2;
  int required = connected ? n - 1 : 0;
  assert(required <= m && m <= total);
  std::vector<std::pair<int, int>> edges;
  if (connected) {
    edges = rand_tree(n);
  }
  std::unordered_set<int64_t> required_edges;
  required_edges.reserve(required);
  for (auto &[u, v] : edges) {
    if ((directed && rand_int(0, 1)) || (!directed && u > v)) {
      std::swap(u, v);
    }
    required_edges.insert(static_cast<int64_t>(u) * n + v);
  }
  int64_t available = total - required;
  int extra = m - required;
  bool complement = extra > available / 2;
  int sample = complement ? static_cast<int>(total - m) : extra;
  // For dense graphs, sample the omitted non-tree edges instead of the included ones.
  std::unordered_set<int64_t> sampled;
  sampled.reserve(sample);
  while (static_cast<int>(sampled.size()) < sample) {
    int u = rand_int(0, n - 1), v = rand_int(0, n - 2);
    if (v >= u) {
      v++;
    }
    if (!directed && u > v) {
      std::swap(u, v);
    }
    int64_t key = static_cast<int64_t>(u) * n + v;
    if (!required_edges.count(key)) {
      sampled.insert(key);
    }
  }
  edges.reserve(m);
  if (!complement) {
    for (int64_t key : sampled) {
      edges.emplace_back(static_cast<int>(key / n), static_cast<int>(key % n));
    }
  } else {
    edges.clear();
    for (int u = 0; u < n; u++) {
      for (int v = directed ? 0 : u + 1; v < n; v++) {
        if (u != v && !sampled.count(static_cast<int64_t>(u) * n + v)) {
          edges.emplace_back(u, v);
        }
      }
    }
  }
  std::shuffle(edges.begin(), edges.end(), rng);
  return edges;
}

std::vector<int> rand_composition(int parts, int sum, int min_part) {
  int64_t remaining = static_cast<int64_t>(sum) - static_cast<int64_t>(parts) * min_part;
  assert(parts >= 1 && min_part >= 0 && remaining >= 0 && remaining + parts - 2 <= INT_MAX);
  if (parts == 1) {
    return {sum};
  }
  auto cuts = rand_distinct(parts - 1, 0, static_cast<int>(remaining) + parts - 2);
  std::sort(cuts.begin(), cuts.end());
  std::vector<int> values;
  int prev = -1;
  for (int cut : cuts) {
    values.push_back(min_part + cut - prev - 1);
    prev = cut;
  }
  values.push_back(min_part + static_cast<int>(remaining) + parts - 2 - prev);
  return values;
}
\end{lstlisting}
A stress test compares an optimized implementation against a trusted brute force implementation on many generated inputs. Small tests make brute force practical and produce manageable counterexamples.

\begin{compactitem}
  \item \inlinecode{stress(trials, gen, solve, brute, equal = std::equal\_to\textless{}\textgreater{}())} runs \inlinecode{trials} tests, calling \inlinecode{gen(trial)} with each 1-based trial number and comparing \inlinecode{solve(test)} with \inlinecode{brute(test)} using \inlinecode{equal}. It returns a \inlinecode{StressFailure} containing the number, first failing test, and actual and expected results. The generated test type must be copyable so each solution receives an independent input, and \inlinecode{equal} must accept the respective results of \inlinecode{solve(test)} and \inlinecode{brute(test)}. If all generated tests pass, it returns \inlinecode{std::nullopt}.
\end{compactitem}

\textbf{Time Complexity}: $O(t \cdot (G + X + S + B + C))$ per call in the worst case, where $t$ is \inlinecode{trials}, $X$ is the cost of copying a generated test, and $G$, $S$, $B$, and $C$ are the costs of \inlinecode{gen}, \inlinecode{solve}, \inlinecode{brute}, and \inlinecode{equal}, respectively.

\textbf{Space Complexity}: $O(x + a + e)$ for the test copies and results, where $x$, $a$, and $e$ are their respective sizes, in addition to the working space used by the callbacks.

\begin{lstlisting}[style=cpp]
template<typename Test, typename Actual, typename Expected>
struct StressFailure {
  int trial;
  Test test;
  Actual actual;
  Expected expected;
};

template<typename Generate, typename Solve, typename Brute, typename Equal = std::equal_to<>>
auto stress(int trials, Generate gen, Solve solve, Brute brute, Equal equal = {}) {
  assert(trials >= 0);
  using Test = std::decay_t<decltype(gen(1))>;
  using Actual = std::decay_t<decltype(solve(std::declval<Test &>()))>;
  using Expected = std::decay_t<decltype(brute(std::declval<Test &>()))>;
  using Failure = StressFailure<Test, Actual, Expected>;
  for (int trial = 1; trial <= trials; ++trial) {
    auto test = gen(trial);
    auto solve_test = test, brute_test = test;
    auto actual = solve(solve_test);
    auto expected = brute(brute_test);
    if (!equal(actual, expected)) {
      return std::optional<Failure>{
          Failure{trial, std::move(test), std::move(actual), std::move(expected)}
      };
    }
  }
  return std::optional<Failure>{};
}
\end{lstlisting}
\Needspace*{4\baselineskip}
\textbf{Example Usage}\par
\begin{lstlisting}[style=cpp]
using namespace std;

bool EQ(double a, double b) {
  return a == b || fabs(a - b) <= 1e-9;
}

int main() {
  rng.seed(1234567);  // Fixed seed for reproducibility.

  double real = rand_real(-1.0, 1.0);
  assert(-1.0 <= real && real < 1.0);
  int biased = rand_biased(1, 100000, 3);
  assert(1 <= biased && biased <= 100000);
  assert(rand_weighted({20, 20, 30}) < 3);
  auto values = rand_vec(8, -5, 5);
  assert(values.size() == 8 && all_of(values.begin(), values.end(), [](int x) {
           return -5 <= x && x <= 5;
         }));
  assert(rand_str(12, "abc").size() == 12);
  auto graph = rand_graph(5, 7);
  assert(graph.size() == 7 && all_of(graph.begin(), graph.end(), [](auto edge) {
           return edge.first < edge.second;
         }));
  assert(rand_tree(6).size() == 5);
  assert(rand_graph(6, 8, false, true).size() == 8);
  auto composition = rand_composition(4, 20, 2);
  assert(accumulate(composition.begin(), composition.end(), 0) == 20);
  assert(*min_element(composition.begin(), composition.end()) >= 2);

  auto generate = [](int trial) {
    int n = rand_int(1, 8), kind = trial <= 3 ? trial - 1 : rand_int(0, 2);
    switch (kind) {
      case 0:
        return rand_perm(n, 1);
      case 1:
        return rand_distinct(n, -5, 5);
      default:
        break;
    }
    static const vector<int> boundary_values{-5, 0, 5};
    bernoulli_distribution use_boundary(0.25);
    vector<int> a(n);
    for (int &x : a) {
      x = use_boundary(rng) ? rand_choice(boundary_values) : rand_int(-5, 5);
    }
    return a;
  };
  auto solve = [](vector<int> a) {
    sort(a.begin(), a.end());
    return a;
  };
  auto brute = [](vector<int> a) {
    do {
      if (is_sorted(a.begin(), a.end())) {
        return a;
      }
    } while (next_permutation(a.begin(), a.end()));
    sort(a.begin(), a.end());
    return a;
  };

  assert(!stress(100, generate, solve, brute));
  assert(!stress(
      1, [](int) { return 0; }, [](int) { return 1.0; }, [](int) { return 1.0 + 1e-12; }, EQ
  ));

  // Mutating solvers still receive independent copies of the generated test.
  auto mutate = [](vector<int> &a) { return ++a[0]; };
  auto unchanged = [](const vector<int> &a) { return a[0]; };
  auto failure = stress(1, [](int) { return vector<int>{1}; }, mutate, unchanged);
  assert(failure && failure->trial == 1 && failure->test == vector<int>{1});
  assert(failure->actual == 2 && failure->expected == 1);
  return 0;
}
\end{lstlisting}

\section{\texorpdfstring{Benchmarking}{8.5 Benchmarking}}
\setcounter{section}{5}
\label{sec:8.5}
Small timing helpers for local profiling and codebook examples. Use \inlinecode{std::chrono::steady\_clock} instead of wall-clock time so elapsed durations are monotonic.

\begin{compactitem}
  \item \inlinecode{Timer()} starts a timer immediately.
  \item \inlinecode{timer.reset()} restarts the timer.
  \item \inlinecode{timer.elapsed()} returns the elapsed time in milliseconds as a double.
  \item \inlinecode{benchmark(iterations, f)} runs \inlinecode{f()} \inlinecode{iterations} times and returns the average elapsed milliseconds per call. \inlinecode{iterations} must be positive.
\end{compactitem}

Benchmarks should be treated as local measurements, not proofs of asymptotic performance. Warm-up, CPU frequency scaling, compiler optimization, and input choice can dominate tiny timings.

\begin{lstlisting}[style=cpp]
#include <cassert>
#include <chrono>

class Timer {
  using Clock = std::chrono::steady_clock;
  Clock::time_point start_time;

 public:
  Timer() { reset(); }

  void reset() { start_time = Clock::now(); }

  double elapsed() const {
    return std::chrono::duration<double, std::milli>(Clock::now() - start_time).count();
  }
};

template<typename Fn>
double benchmark(int iterations, Fn f) {
  assert(iterations > 0);
  Timer timer;
  for (int i = 0; i < iterations; ++i) {
    f();
  }
  return timer.elapsed() / iterations;
}
\end{lstlisting}
\Needspace*{4\baselineskip}
\textbf{Example Usage}\par
\begin{lstlisting}[style=cpp]
#include <cassert>
#include <numeric>
#include <vector>
using namespace std;

int main() {
  Timer timer;
  vector<int> v(1000);
  iota(v.begin(), v.end(), 0);
  assert(timer.elapsed() >= 0.0);

  volatile long long sink = 0;
  double avg_time = benchmark(5, [&] { sink += accumulate(v.begin(), v.end(), 0LL); });
  assert(avg_time >= 0.0);
  assert(sink > 0);
  return 0;
}
\end{lstlisting}

\section{\texorpdfstring{GNU Policy-Based Data Structures}{8.6 GNU Policy-Based Data Structures}}
\setcounter{section}{6}
\label{sec:8.6}
GNU policy-based data structures (PBDS) provide non-standard associative containers and priority queues, including order-statistic trees and fast hash tables. They are not part of the C++ standard library, but they are available on many GNU C++ contest judges.

\begin{compactitem}
  \item \inlinecode{HashMap\textless{}K, V\textgreater{}} and \inlinecode{HashSet\textless{}K\textgreater{}} behave like faster, non-standard alternatives to \inlinecode{std::unordered\_map\textless{}K, V\textgreater{}} and \inlinecode{std::unordered\_set\textless{}K\textgreater{}}, using \inlinecode{gp\_hash\_table}. The default \inlinecode{SplitMix64Hasher} is for integer-like keys; pass a custom hash for other key types. It randomizes the hash seed to avoid weak adversarial integer hashes.
  \item \inlinecode{OrderedSet\textless{}T\textgreater{}} and \inlinecode{OrderedMap\textless{}K, V\textgreater{}} behave like \inlinecode{std::set\textless{}T\textgreater{}} and \inlinecode{std::map\textless{}K, V\textgreater{}} but also support \inlinecode{find\_by\_order(k)} and \inlinecode{order\_of\_key(x)}.
  \item \inlinecode{OrderedMultiset\textless{}T\textgreater{}} behaves like \inlinecode{std::multiset\textless{}T\textgreater{}} by storing duplicates as (value, unique ID) pairs.
  \item \inlinecode{OrderedMultimap\textless{}K, V\textgreater{}} behaves like \inlinecode{std::multimap\textless{}K, V\textgreater{}} by storing duplicate keys as (key, unique ID) pairs.
  \item \inlinecode{PairingHeap\textless{}T\textgreater{}} is a meldable priority queue; \inlinecode{push()} returns an iterator that can be passed to \inlinecode{modify()} or \inlinecode{erase()}.
  \item \inlinecode{BinaryHeap\textless{}T\textgreater{}}, \inlinecode{BinomialHeap\textless{}T\textgreater{}}, and \inlinecode{RcBinomialHeap\textless{}T\textgreater{}} are alternative GNU priority queue tags with the same interface.
\end{compactitem}

The order-statistic operations use GNU PBDS conventions: \inlinecode{find\_by\_order(k)} is 0-based and \inlinecode{order\_of\_key(x)} returns the number of keys strictly less than \inlinecode{x}, even if \inlinecode{x} is absent. Raw PBDS ordered trees return iterators from \inlinecode{find\_by\_order(k)}; the multiset and multimap wrappers below unwrap that iterator into a value. The priority queue \inlinecode{join()} operation melds another heap into the current one. Since this depends on \inlinecode{\_\_gnu\_pbds}, keep a standard library fallback in mind when portability matters.

\Needspace*{3\baselineskip}
\textbf{Time Complexity}:
\begin{compactitem}
  \item \inlinecode{HashMap}/\inlinecode{HashSet}: same expected bounds as \inlinecode{std::unordered\_map}/\inlinecode{std::unordered\_set} (expected $O(1)$, worst-case $O(n)$).
  \item \inlinecode{OrderedSet}/\inlinecode{OrderedMap}: same as \inlinecode{std::set}/\inlinecode{std::map}, with $O(\log n)$ for \inlinecode{find\_by\_order} and \inlinecode{order\_of\_key}.
  \item \inlinecode{OrderedMultiset}/\inlinecode{OrderedMultimap}: same as \inlinecode{std::multiset}/\inlinecode{std::multimap}, with $O(\log n)$ for \inlinecode{find\_by\_order} and \inlinecode{order\_of\_key}.
  \item $O(1)$ amortized per call to \inlinecode{push()} and \inlinecode{join()} for \inlinecode{PairingHeap}.
  \item $O(\log n)$ amortized per call to \inlinecode{pop()}, \inlinecode{modify()}, and \inlinecode{erase()} for \inlinecode{PairingHeap}.
\end{compactitem}

\textbf{Space Complexity}: $O(n)$ for all containers.

\begin{lstlisting}[style=cpp]
#include <cassert>
#include <chrono>
#include <cstddef>
#include <cstdint>
#include <ext/pb_ds/assoc_container.hpp>
#include <ext/pb_ds/priority_queue.hpp>
#include <ext/pb_ds/tree_policy.hpp>
#include <functional>
#include <utility>

struct SplitMix64Hasher {
  static uint64_t splitmix64(uint64_t x) {
    x += 0x9e3779b97f4a7c15ULL;
    x = (x ^ (x >> 30)) * 0xbf58476d1ce4e5b9ULL;
    x = (x ^ (x >> 27)) * 0x94d049bb133111ebULL;
    return x ^ (x >> 31);
  }

  std::size_t operator()(uint64_t x) const {
    static const uint64_t RAND_SEED = std::chrono::steady_clock::now().time_since_epoch().count();
    return splitmix64(x + RAND_SEED);
  }
};

template<typename K, typename Hash = SplitMix64Hasher>
using HashSet = __gnu_pbds::gp_hash_table<K, __gnu_pbds::null_type, Hash>;

template<typename K, typename V, typename Hash = SplitMix64Hasher>
using HashMap = __gnu_pbds::gp_hash_table<K, V, Hash>;

template<typename K, typename V, typename Compare = std::less<K>>
using OrderedMap = __gnu_pbds::tree<
    K, V, Compare, __gnu_pbds::rb_tree_tag, __gnu_pbds::tree_order_statistics_node_update>;

template<typename T, typename Compare = std::less<T>>
using OrderedSet = OrderedMap<T, __gnu_pbds::null_type, Compare>;

template<typename T>
class OrderedMultiset {
  OrderedSet<std::pair<T, int>> tree;
  int timer = 0;

 public:
  void insert(const T &x) { tree.insert({x, timer++}); }

  bool erase_one(const T &x) {
    auto it = tree.lower_bound({x, -1});
    if (it == tree.end() || it->first != x) {
      return false;
    }
    tree.erase(it);
    return true;
  }

  T find_by_order(int k) const {
    assert(0 <= k && k < size());
    return tree.find_by_order(k)->first;
  }

  int order_of_key(const T &x) const { return tree.order_of_key({x, -1}); }
  int size() const { return static_cast<int>(tree.size()); }
};

template<typename K, typename V>
class OrderedMultimap {
  OrderedMap<std::pair<K, int>, V> tree;
  int timer = 0;

 public:
  void insert(const K &key, const V &value) { tree.insert({{key, timer++}, value}); }

  bool erase_one(const K &key) {
    auto it = tree.lower_bound({key, -1});
    if (it == tree.end() || it->first.first != key) {
      return false;
    }
    tree.erase(it);
    return true;
  }

  std::pair<K, V> find_by_order(int k) const {
    assert(0 <= k && k < size());
    auto it = tree.find_by_order(k);
    return {it->first.first, it->second};
  }

  int order_of_key(const K &x) const { return tree.order_of_key({x, -1}); }
  int size() const { return static_cast<int>(tree.size()); }
};

template<typename T, typename Compare = std::less<T>>
using PairingHeap = __gnu_pbds::priority_queue<T, Compare, __gnu_pbds::pairing_heap_tag>;

template<typename T, typename Compare = std::less<T>>
using BinaryHeap = __gnu_pbds::priority_queue<T, Compare, __gnu_pbds::binary_heap_tag>;

template<typename T, typename Compare = std::less<T>>
using BinomialHeap = __gnu_pbds::priority_queue<T, Compare, __gnu_pbds::binomial_heap_tag>;

template<typename T, typename Compare = std::less<T>>
using RcBinomialHeap = __gnu_pbds::priority_queue<T, Compare, __gnu_pbds::rc_binomial_heap_tag>;
\end{lstlisting}
\Needspace*{4\baselineskip}
\textbf{Example Usage}\par
\begin{lstlisting}[style=cpp]
int main() {
  OrderedSet<int> s;
  s.insert(10);
  s.insert(20);
  s.insert(30);
  s.insert(20);  // Duplicate ignored, like std::set.
  // Net new over std::set: order-statistic queries.
  assert(s.size() == 3);
  assert(*s.find_by_order(0) == 10);
  assert(*s.find_by_order(1) == 20);
  assert(s.order_of_key(25) == 2);
  assert(s.order_of_key(100) == 3);

  HashMap<int, int> hm;
  hm[10] = 1;
  hm[20] = 2;
  assert(hm[10] == 1);
  HashSet<int> hs;
  hs.insert(7);
  assert(hs.find(7) != hs.end());

  OrderedMultiset<int> ms;
  ms.insert(5);
  ms.insert(5);
  ms.insert(7);
  assert(ms.size() == 3);
  // Net new over std::multiset: order statistics with duplicates.
  assert(ms.order_of_key(7) == 2);
  assert(ms.find_by_order(1) == 5);
  assert(ms.erase_one(5));
  assert(ms.order_of_key(7) == 1);
  assert(ms.erase_one(5));
  assert(!ms.erase_one(5));
  assert(ms.find_by_order(0) == 7);

  OrderedMultimap<int, char> mp;
  mp.insert(5, 'a');
  mp.insert(5, 'b');
  mp.insert(7, 'c');
  // Net new over std::multimap: order statistics with duplicate keys.
  assert(mp.order_of_key(7) == 2);
  assert(mp.find_by_order(1).first == 5);
  assert(mp.erase_one(5));
  assert(mp.order_of_key(7) == 1);
  assert(mp.erase_one(5));
  assert(!mp.erase_one(5));
  assert(mp.find_by_order(0).first == 7);

  PairingHeap<int> h1, h2;
  auto it = h1.push(10);
  auto erased = h1.push(5);
  h1.modify(it, 20);  // Net new over std::priority_queue: update through a handle.
  h1.erase(erased);   // Also net new: remove an arbitrary heap item through a handle.
  h2.push(7);
  h1.join(h2);  // Also net new: meld another heap into this one.
  assert(h2.empty());
  assert(h1.top() == 20);
  h1.pop();
  assert(h1.top() == 7);
  return 0;
}
\end{lstlisting}

\section{\texorpdfstring{Multithreading}{8.7 Multithreading}}
\setcounter{section}{7}
\label{sec:8.7}
Basic multithreading can speed up independent local work such as stress tests, batch experiments, or offline preprocessing. It is rarely useful in standard contest submissions, and many judges restrict threads, so use it only when the environment allows it.

\begin{compactitem}
  \item \inlinecode{parallel\_for(n, threads, f)} runs \inlinecode{f(i)} once for each $0 \leq {\inlinecodemath{i}} < {\inlinecodemath{n}}$, splitting the indices among at most \inlinecode{threads} worker threads.
  \item \inlinecode{parallel\_cases(cases, threads, solve)} solves independent test cases in parallel and returns the outputs in the original case order, using at most \inlinecode{threads} workers and calling \inlinecode{solve(case, output)} for each case.
\end{compactitem}

Each call to \inlinecode{f(i)} must be safe to run concurrently with the others. If it writes shared state, use separate output slots or synchronization.

\begin{lstlisting}[style=cpp]
#include <algorithm>
#include <sstream>
#include <string>
#include <thread>
#include <vector>

template<typename Fn>
void parallel_for(int n, int threads, Fn f) {
  threads = std::max(1, std::min(threads, n));
  std::vector<std::thread> workers;
  workers.reserve(threads);
  for (int id = 0; id < threads; ++id) {
    workers.emplace_back([=, &f] {
      for (int i = id; i < n; i += threads) {
        f(i);
      }
    });
  }
  for (auto &worker : workers) {
    worker.join();
  }
}

template<typename Case, typename Solve>
std::vector<std::string> parallel_cases(const std::vector<Case> &cases, int threads, Solve solve) {
  std::vector<std::string> out(cases.size());
  parallel_for(static_cast<int>(cases.size()), threads, [&](int i) {
    std::ostringstream os;
    solve(cases[i], os);
    out[i] = os.str();
  });
  return out;
}
\end{lstlisting}
\Needspace*{4\baselineskip}
\textbf{Example Usage}\par
\begin{lstlisting}[style=cpp]
#include <cassert>
#include <numeric>
using namespace std;

int main() {
  vector<long long> squares(1000);
  parallel_for(static_cast<int>(squares.size()), 4, [&](int i) { squares[i] = 1LL * i * i; });
  assert(squares[17] == 289);
  assert(accumulate(squares.begin(), squares.end(), 0LL) == 332833500);

  vector<int> cases{3, 1, 4, 2};
  auto out =
      parallel_cases(cases, 2, [](int x, ostringstream &os) { os << x << "^2 = " << x * x; });
  assert((out == vector<string>{"3^2 = 9", "1^2 = 1", "4^2 = 16", "2^2 = 4"}));
  return 0;
}
\end{lstlisting}

\section{\texorpdfstring{Bump Allocation}{8.8 Bump Allocation}}
\setcounter{section}{8}
\label{sec:8.8}
A typed bump pool replaces many small heap allocations with consecutive slots in one fixed vector. This simple form is useful for pointer-based trees, tries, and linked structures whose nodes all live until the algorithm finishes. Allocation advances one index, and individual nodes are never freed.

\begin{compactitem}
  \item \inlinecode{make\_node(value, next = nullptr)} returns a pointer to a new node initialized with \inlinecode{value} and \inlinecode{next}.
  \item \inlinecode{BumpArena\textless{}N\textgreater{}} owns an inline $N$-byte monotonic buffer. Pass \inlinecode{arena.resource()} to a \inlinecode{std::pmr} container such as \inlinecode{std::pmr::vector\textless{}T\textgreater{}} or \inlinecode{std::pmr::map\textless{}K, V\textgreater{}}.
\end{compactitem}

The vector is created at its final size and never resized, so returned pointers remain valid. Adjust \inlinecode{MAX\_NODES} and the fields of \inlinecode{Node} for the problem at hand.

For standard containers, the C++17 polymorphic allocator interface avoids both a global \inlinecode{operator new} override and the full allocator protocol. An arena must outlive every container using it; declare a large arena statically if it would exceed the stack limit. The wrapper uses \inlinecode{std::pmr::null\_memory\_resource()} as its upstream resource, so exhausting the buffer throws \inlinecode{std::bad\_alloc} instead of silently allocating from the heap.

\Needspace*{3\baselineskip}
\textbf{Time Complexity}:
\begin{compactitem}
  \item $O({\inlinecodemath{MAX\_NODES}})$ for construction of the pool.
  \item $O(1)$ per call to \inlinecode{make\_node()}.
\end{compactitem}

\Needspace*{3\baselineskip}
\textbf{Space Complexity}:
\begin{compactitem}
  \item $O({\inlinecodemath{MAX\_NODES}})$ for storage of the pool.
  \item $O(N)$ for storage of each \inlinecode{BumpArena\textless{}N\textgreater{}}.
\end{compactitem}

\begin{lstlisting}[style=cpp]
#include <cassert>
#include <cstddef>
#include <memory_resource>
#include <vector>

struct Node {
  int value;
  Node *next;
};

static constexpr int MAX_NODES = 100000;
static std::vector<Node> node_pool(MAX_NODES);
static int node_count = 0;

Node *make_node(int value, Node *next = nullptr) {
  assert(node_count < MAX_NODES);
  node_pool[node_count] = {value, next};
  return &node_pool[node_count++];
}

template<std::size_t N>
class BumpArena {
  alignas(std::max_align_t) std::byte buffer[N];
  std::pmr::monotonic_buffer_resource arena{buffer, N, std::pmr::null_memory_resource()};

 public:
  std::pmr::memory_resource *resource() { return &arena; }
};
\end{lstlisting}
\Needspace*{4\baselineskip}
\textbf{Example Usage}\par
\begin{lstlisting}[style=cpp]
#include <map>
using namespace std;

int main() {
  Node *tail = make_node(2);
  Node *head = make_node(1, tail);
  assert(head->value == 1 && head->next->value == 2);

  BumpArena<1024> arena;
  pmr::vector<int> values(arena.resource());
  values.reserve(4);
  values.insert(values.end(), {1, 2, 3, 4});
  assert(values.size() == 4 && values.back() == 4);

  pmr::map<int, int> counts(arena.resource());
  counts.emplace(7, 2);
  assert(counts[7] == 2);
  return 0;
}
\end{lstlisting}
